%! Projections onto Subspaces — Unit 4, Lesson 4.2 — Student Packet Cover Sheet
\documentclass[10pt]{article}
\usepackage{linalg-article}
\usepackage{linalg-boxes}
\usepackage{ltablex}
\keepXColumns

\begin{document}

% Full-bleed burgundy banner
\begin{tikzpicture}[remember picture, overlay]
  \fill[burgundy] (current page.north west)
    rectangle ([yshift=-0.9in]current page.north east);
\end{tikzpicture}
\vspace{-0.8in}

\begin{center}
  {\color{white}\textbf{\LARGE Linear Algebra}}\\[4pt]
  {\color{blushmid}\large\textbf{Unit 4 \quad Orthogonality}}\\[6pt]
  {\color{white}\textbf{\large Lesson 4.2 \quad Projections onto Subspaces}}
\end{center}

\vspace{0.22in}
\namedateperiod
\vspace{0.18in}

\begin{learningtargetbox}
\small
\begin{enumerate}[leftmargin=1.5em, itemsep=5pt]
  \item \ldots{} find the \textbf{projection} of a vector onto a \textbf{line} using $\hat{x}=\frac{\mathbf{a}\cdot\mathbf{b}}{\mathbf{a}\cdot\mathbf{a}}$, and locate the perpendicular \textbf{error}.
  \item \ldots{} project onto a \textbf{subspace} by setting up and solving the \textbf{normal equations} $A^{\mathsf{T}}A\hat{\mathbf{x}}=A^{\mathsf{T}}\mathbf{b}$.
  \item \ldots{} explain \emph{why} the closest point is the one whose error is \textbf{perpendicular} to the space.
\end{enumerate}
\end{learningtargetbox}

\vspace{0.14in}

\begin{tocbox}
\small
\renewcommand{\arraystretch}{1.5}
\begin{tabularx}{\linewidth}{@{} c l X r @{}}
\rowcolor{blush}
\textbf{\#} & \textbf{Component} & \textbf{Description} & \textbf{Score} \\
\midrule
1 & Warm-Up        & Spiral review: perpendicular test, two key dot products, solving a $2\times2$ system & \blank{1.2cm} \\
2 & Guided Notes   & Closest point, projection onto a line \& subspace, the error, normal equations & \blank{1.2cm} \\
3 & Group Activity & Drop the Perpendicular --- projecting onto a line, a matrix, and a plane & \blank{1.2cm} \\
4 & Exit Ticket    & Quick check on a line projection, the error's meaning, and normal equations & \blank{1.2cm} \\
5 & Homework       & Line \& subspace projections, projection matrices, and where the error lives & \blank{1.2cm} \\
\midrule
  & \multicolumn{2}{r}{\textbf{Total}} & \blank{1.2cm} \\
\end{tabularx}
\end{tocbox}

\vspace{0.10in}

\begin{remindbox}
\small One idea drives the lesson: the \textbf{closest} point $\mathbf{p}$ in a line or subspace is the
one whose leftover \textbf{error} $\mathbf{e}=\mathbf{b}-\mathbf{p}$ is \emph{perpendicular} to the space.
Turning that right angle into an equation gives $\hat{x}=\frac{\mathbf{a}\cdot\mathbf{b}}{\mathbf{a}\cdot\mathbf{a}}$
for a line and the \textbf{normal equations} $A^{\mathsf{T}}A\hat{\mathbf{x}}=A^{\mathsf{T}}\mathbf{b}$ for
a subspace --- and the error always lands in the orthogonal complement (Lesson~4.1).
\end{remindbox}

\end{document}
